% Source: https://artofproblemsolving.com/community/c3h2577722_ajax31_birthday_problems
% Stable identity: aops-topic-2577722-q04
Place the half-disc as $\{(x,y):\ y\ge0,\ x^2+y^2\le25\}$. Every point of it has $0\le y\le5$.

\emph{Upper bound.} Let an equilateral triangle of side $s$ lie in the half-disc. Its three sides point in three directions that are $60^\circ$ apart, so they split the $180^\circ$ of possible directions into three arcs of $60^\circ$. Hence some side makes an angle of at most $30^\circ$ with the vertical, and so an angle of at least $60^\circ$ with the horizontal. The two endpoints of that side differ in height by at least $s\sin60^\circ=\dfrac{\sqrt3}2s$. Both endpoints lie in the half-disc, so this difference is at most $5$. Hence $s\le\dfrac{10}{\sqrt3}$.

\emph{Construction.} The triangle $PQT$ with $P=\left(-\dfrac5{\sqrt3},0\right)$, $Q=\left(\dfrac5{\sqrt3},0\right)$, $T=(0,5)$ has side $\dfrac{10}{\sqrt3}$ and height $5$, and lies in the half-disc: $T$ is on the arc and $P$, $Q$ are on the diameter.

So the largest area is
\[\dfrac{\sqrt3}4\cdot\left(\dfrac{10}{\sqrt3}\right)^2=\boxed{\dfrac{25\sqrt3}{3}}\approx14.43.\]
(The dashed triangle in the sketch is a tilted triangle; any such triangle is no larger.)

\begin{asy}
usepackage("amsmath");
size(5.5cm);
defaultpen(fontsize(10pt));
real R=5;pair O=(0,0);
real s=10/sqrt(3);
draw(arc(O,R,0,180)--cycle,linewidth(0.8));
pair P=(-s/2,0),Q=(s/2,0),T=(0,R);
fill(P--Q--T--cycle,gray(0.88));draw(P--Q--T--cycle,linewidth(0.8));
pair A1=(-2.6,0.6),A2=A1+3.8*dir(20),A3=A1+3.8*dir(80);
draw(A1--A2--A3--cycle,dashed);
draw((5.8,0)--(5.8,5),Arrows(5));label("$5$",(5.8,2.5),E);
draw((5.6,5)--(6,5));draw((5.6,0)--(6,0));
dot(O);label("$O$",O,S);label("$P$",P,S);label("$Q$",Q,S);label("$T$",T,N);
\end{asy}
